\scanpage{37}
\mspage{35}
% archive sheet number: 36
8. \uline{Théorème de monodromie : démonstration arithmétique}

\uline{Th.} 8.1. Supposons que $V$ soit limite inductive \inserted{filtrante} de sous-anneaux locaux $V_i$, \deleted{l'idéal} \inserted{d'idéal maximal $\unsure{\mathfrak m_i|m_i}=\mathfrak m\cap V_i$} contenant l'uniformisante $t$, tels que les deux conditions suivantes soient satisfaites :

a) \deleted{\unsure{le div}} $V_i$ \unsure{est} \inserted{et $V_i/tV_i$ sont} réguliers, \deleted{et le diviseur $\operatorname{div}(t)$ \unsure{est} régulier, i.e. $t$ fait partie d'un système de param. de $V_i$.} \deleted{dans $\operatorname{Spec}V_i$, est à croisements normaux.}

b) Le corps \inserted{résiduel} $k_i=k(V_i)$ est tel que aucune extension finie de $k_i$ ne contienne toutes les racines $\ell^{\unsure{\nu|n}}$-ièmes de l'unité, pour $\forall$ \unsure{$\ell$ entier $>1$}.

Alors le th. de monodromie \unsure{sous} la forme 7.1. est vrai pour des schémas \deleted{de t.f.} \inserted{\uline{\unsure{propre|propres}}} $X$ sur $K$, \deleted{tel que $\unsure{H^i|H^0}(X_{\overline K},\unsure{\mathbb Z/\ell\mathbb Z|Z/\ell Z})$ est fini} \inserted{lorsque l'opération de $I$ sur $H^i(X_{\overline K},\mathbb Z_\ell)\otimes_{\mathbb Z_\ell}\unsure{\mathbb Q_\ell|Q_\ell}$ se fait à travers $I_t$,} \deleted{dans} \uline{ou} que les $V_\alpha$ sont des anneaux de val. discrète.

\uline{Corollaire} 8.2. Le théorème de monodromie \inserted{sous la forme 7.1.} est vrai \deleted{lorsque $H^i(X_{\overline K},\mathbb Z/\ell\mathbb Z)$ est fini,} \inserted{pour $X_K$ propre sur $K$} dans chacun des cas suivants :

a) $V$ est \inserted{le hensélisé ou le hensélisé strict de} l'anneau local d'un point d'un schéma $S$ de type fini sur un corps $k$, ou sur \deleted{\unsure{Spec}} \inserted{un localisé de} $\mathbb Z$.

b) $V$ est de car. nulle, (i.e. \inserted{$k$ est} de caractéristique \deleted{\unsure{résiduelle}} nulle).

\deleted{Cas a)}
\unsure{Démonstration}. 8.2.

\uline{Cas} a). \deleted{Supposons $V$ le hensélisé de l'anneau local d'un \unsure{point}. \unsure{Si} $S$ est de type fini sur $\operatorname{Spec}\mathbb Z$,}
\unsure{Supposons} d'abord que l'anneau de val. \deleted{\unsure{discr.}} $V$ soit l'anneau local d'un schéma $S$ de type fini sur \deleted{\unsure{Spec}} \inserted{un localisé} $\mathbb Z$, ou sur un corps $k_0$. Dans le premier cas, \inserted{ou si $k_0$ est de type fini sur le corps premier,} le corps résiduel $k(V)$ est de type fini sur le corps premier, donc $V$ satisfait la condition
